\chapter{Chaînes de Markov et distributions invariantes}
\section*{Exercice 35.1 -- Convention de transition}
[Base] Pour $P=\begin{pmatrix}0{,}8&0{,}2\\0{,}3&0{,}7\end{pmatrix}$ et la distribution ligne $\mu_0=(1,0)$, calculer $\mu_1,\mu_2$.

\section*{Exercice 35.2 -- Distribution invariante}
[Méthode] Pour la matrice précédente, déterminer la distribution invariante et la limite de $\mu_n$.

\section*{Exercice 35.3 -- Périodicité}
[Critique] Pour $P=\begin{pmatrix}0&1\\1&0\end{pmatrix}$, trouver une distribution invariante et expliquer pourquoi la distribution issue de $(1,0)$ ne converge pas.

\section*{Exercice 35.4 -- Trois états}
[Méthode] Pour $P=\begin{pmatrix}1/2&1/2&0\\0&1/2&1/2\\1/2&0&1/2\end{pmatrix}$, vérifier que $(1/3,1/3,1/3)$ est invariante et calculer la distribution après deux pas depuis l'état 1.

\section*{Exercice 35.5 -- Probabilité d’un trajet}
[Application] Pour la matrice de l'exercice 35.1, partir de l'état 1 et calculer la probabilité du trajet $1,2,2,1$.

\section*{Exercice 35.6 -- État absorbant}
[Approfondissement] Pour $P=\begin{pmatrix}q&1-q\\0&1\end{pmatrix}$, $0<q<1$, partir du premier état. Calculer la probabilité d'y rester après $n$ pas et le temps moyen avant absorption.
